\input python_2_stepper_style.tex

\section*{Introduction}

The \emph{stepper} shows how a Python \S 1 or Python \S 2 program is evaluated
by the \emph{substitution model}, as used in Chapters 1 and 2 of SICPy: the
program itself is rewritten, one small step at a time, until it cannot be
rewritten any further. In contrast to the CSE machine, no environments exist:
a name is replaced by its value, and a function application is replaced by
the function's body in which the arguments have been substituted for the
parameters. The substitution model does not support assignment to names that
have already been used; that is the reason for Python \S 3 to move on to the
environment model, visualized by the CSE machine and the environment stepper.

This document specifies the rewriting that the stepper performs. It follows
the specification of the Source \S 2 stepper, with the differences that arise
from Python's syntax: programs are sequences of statements without semicolons,
function bodies are suites that end in \sCode{return}, and a statement that
consists of a value has no result.

\section*{Reduction}

The \emph{reducer} $\Rightarrow$ is a partial function from
programs/statements/expressions to
programs/statements/expressions (with slight abuse of notation
via overloading)
and $\Rightarrow^*$ is its reflexive transitive closure.
A \emph{reduction} is a sequence of programs
$p_1 \Rightarrow \cdots \Rightarrow p_n$,
where $p_n$ is not reducible, i.e. there is no
program $q$ such that $p_n \Rightarrow q$.
Here, the program $p_n$ is called the \emph{result
of reducing} $p_1$.

If $p_n$, the result of reducing a program, is the empty
sequence of statements, the evaluation is \emph{complete}.
Reduction can get ``stuck'': the result can be a program
that is not empty, which corresponds to a runtime error. In both cases the
stepper ends the reduction with a last step that says
``Evaluation complete'' or ``Evaluation stuck'', respectively.
Unlike in Source, a statement that consists of a value does not
survive as the result of the program: it is removed when it is reached
(see \textbf{Program-discard}); the stepper therefore always ends with
the empty program, unless a runtime error occurs.

A \emph{value} is an integer, float, complex or string literal,
\sCode{True}, \sCode{False}, \sCode{None},
a \sCode{lambda} expression, a function definition
(\sCode{def}), a built-in function, or a pair
\sCode{[}$v_1$\sCode{, }$v_2$\sCode{]} of values. (The stepper shows
the result of \sCode{pair(}$v_1$\sCode{, }$v_2$\sCode{)} in this
notation; \sCode{None} is the empty linked list.) A function that was
declared with a name is shown by that name and reveals its definition
when hovered.

The \emph{substitution} function
$p [ n := v ]$ on programs/statements/expressions
replaces every free occurrence of the name $n$
in $p$ by value $v$. See the section on substitution below. Care must be
taken to introduce and preserve
co-references in this process; substitution can introduce
cyclic references in the result of the substitution. For example,
$n$ may occur free in $v$, in which case
every occurrence of $n$ in $p$
will be replaced by $v$ such that $n$ in $v$ refers cyclically
to the node at which the replacement happens.

Expressions are reduced from left to right, and inner expressions
before the expressions that contain them: in
$e_1 \sCode{ + } e_2$, the expression $e_2$ is only reduced when $e_1$
is a value. Each reduction step is called a \emph{contraction}. The
stepper shows two states for each contraction: first the program
with the expression that is about to be contracted highlighted
(the \emph{redex}), and then the program after the contraction with
the result highlighted.

\subsection*{Programs}

\sEntry{Program-intro}{In a sequence of statements, we can always reduce the
    first statement.}
\gentzen{
    \sVar{statement} \Rightarrow \sVar{statement}'
}{
    \sVar{statement}\ \sVar{statement}_2\ldots \Rightarrow
    \sVar{statement}'\ \sVar{statement}_2\ldots
}

\sEntry{Program-discard}{When the first statement of a program is an
    expression statement whose expression is a value, the statement is
    removed. In Python, a statement does not have a result.}
\gentzen{
    \sVar{v} \mbox{ is a value}
}{
    \sVar{v}\ \sVar{statement}\ldots \Rightarrow
    \sVar{statement}\ldots
}

\sEntry{Eliminate-function-definition}{A function definition as the first
    statement is substituted in the remaining statements, in which its
    name is replaced by the function.}
\gentzen{
             f = \sCode{def } \sVar{name}\sCode{(} x_1 \sCode{, }\ldots\sCode{, }x_n
                 \sCode{):}\ \sVar{suite}
}{
  f\ \sVar{statement} \ldots\ 
  \Rightarrow\ 
  \sVar{statement} \ldots[\sVar{name} := f]
}

\sEntry{Eliminate-assignment}{An assignment of a value as the first
    statement is substituted in the remaining statements.}
\gentzen{
             \sVar{v} \mbox{ is a value}
}{
  \sVar{name}\sCode{ = }\sVar{v}\ \sVar{statement} \ldots\ 
  \Rightarrow \ 
  \sVar{statement} \ldots[\sVar{name} := \sVar{v}]
}

\sEntry{Pass}{A \sCode{pass} statement is removed.}
\gentzen{}{
  \sCode{pass}\ \sVar{statement}\ldots \Rightarrow \sVar{statement}\ldots
}

\sEntry{Import}{An import statement is removed, and the imported names are
  substituted in the remaining statements by the functions of the module.}
\gentzen{}{
  \sCode{from }\sVar{module}\sCode{ import }\sVar{name}\ \sVar{statement}\ldots
  \Rightarrow \sVar{statement}\ldots[\sVar{name} := \sVar{f}]
}

\sEntry{Breakpoint}{The statement \sCode{breakpoint()} is removed. As a
  side-effect, the step before this contraction is marked as a breakpoint, to
  which the \emph{previous breakpoint} and \emph{next breakpoint} buttons of the
  stepper jump. A statement on which the user has set a breakpoint in the
  editor is treated in the same way, in addition to its own contraction: the
  first step that reaches it is marked.}
\gentzen{}{
  \sCode{breakpoint()}\ \sVar{statement}\ldots \Rightarrow \sVar{statement}\ldots
}

\subsection*{Statements: Expression statements}

\sEntry{Expression-statement-reduce}{An expression statement
is reducible if its expression is reducible.}
\gentzen{
  e\ \Rightarrow\ e'
}{
  e \ \Rightarrow \ e'
}

\subsection*{Statements: Assignments}

\sEntry{Evaluate-assignment}{The right-hand expression
in an assignment is evaluated.}
\gentzen{
  \sVar{expression}
  \ \Rightarrow \ 
  \sVar{expression}'
}{
  \sVar{name}\sCode{ = }\sVar{expression}
  \ \Rightarrow \ 
  \sVar{name}\sCode{ = }\sVar{expression}'
}

\subsection*{Statements: Conditionals}

\sEntry{Conditional-statement-predicate}{A conditional statement
is reducible if its predicate is reducible.}
\gentzen{
  e\ \Rightarrow\ e'
}{
  \sCode{if }e\sCode{:}\ \sVar{suite}_1\ \sCode{else:}\ \sVar{suite}_2
  \ \Rightarrow \ 
  \sCode{if }e'\sCode{:}\ \sVar{suite}_1\ \sCode{else:}\ \sVar{suite}_2
}

An \sCode{elif} is the \sCode{else} suite of a conditional statement, which
consists of a conditional statement; an \sCode{if} without
\sCode{else} has an empty alternative suite.

\sEntry{Conditional-statement-consequent}{A conditional statement whose
    predicate is \sCode{True} is replaced by the statements of the consequent
    suite.}
\gentzen{
}{
  \sCode{if True:}\ \sVar{suite}_1\ \sCode{else:}\ \sVar{suite}_2\ \sVar{statement}\ldots
  \ \Rightarrow \ 
  \sVar{suite}_1\ \sVar{statement}\ldots
}

\sEntry{Conditional-statement-alternative}{A conditional statement whose
    predicate is \sCode{False} is replaced by the statements of the
    alternative suite.}
\gentzen{
}{
  \sCode{if False:}\ \sVar{suite}_1\ \sCode{else:}\ \sVar{suite}_2\ \sVar{statement}\ldots
  \ \Rightarrow \ 
  \sVar{suite}_2\ \sVar{statement}\ldots
}

If the predicate is a value other than \sCode{True} and \sCode{False}, the
reduction is stuck with the error \texttt{TypeError: predicate type must be 'bool'}.

\subsection*{Function bodies}

When a function with a suite as its body is applied, the application is
replaced by a \emph{block expression}: the statements of the body, with the
arguments substituted. A block expression is reduced like a program, except
for the following rules.

\sEntry{Block-expression-intro}{A block expression is
reducible if its first statement is reducible. The rules for statements
above apply in the same way as for programs, including
\textbf{Program-discard} for value statements that are not the last statement.}
\gentzen{
  \sVar{statement} \Rightarrow \sVar{statement}'
}{
  \sCode{\{}\ \sVar{statement}\ \sVar{statement}_2\ldots\ \sCode{\}}
  \Rightarrow
  \sCode{\{}\ \sVar{statement}'\ \sVar{statement}_2\ldots\ \sCode{\}}
}

\sEntry{Block-expression-return}{A block expression starting with a
    return statement reduces to the expression of the return statement. The
    statements after the return statement are dropped. A \sCode{return}
    without an expression returns \sCode{None}.}
\gentzen{}{
    \sCode{\{ return } \sVar{e}\ \sVar{statement\ldots}\ \sCode{\}}
    \Rightarrow
    \sVar{e}
}

\sEntry{Block-expression-fall-off}{A block expression that has no statements
   left, or whose last statement is a value that is discarded, reduces to
   \sCode{None}.}
\gentzen{}{
    \sCode{\{ \}} \Rightarrow \sCode{None}
}

As a shortcut, a function whose body starts with a \sCode{return}
statement, as well as every \sCode{lambda} expression, is
applied by replacing the application directly with the returned expression
(or the body of the lambda expression, respectively), without
creating a block expression.

\subsection*{Expressions: Binary operators}

\sEntry{Left-binary-reduce}{An expression with a binary operator
can be reduced if its left sub-expression can be reduced.}
\gentzen{
  e_1 \ \Rightarrow \ e_1'
}{
  e_1\  \sVar{binary-operator} \ e_2
  \ \Rightarrow \ 
  e_1'\  \sVar{binary-operator} \ e_2
}

\sEntry{Right-binary-reduce}{An expression with a binary operator
can be reduced if its left sub-expression is a value and its right
sub-expression can be reduced.}
\gentzen{
  e_2 \ \Rightarrow \ e_2'
}{
  v_1\  \sVar{binary-operator} \ e_2
  \ \Rightarrow \ 
  v_1\  \sVar{binary-operator} \ e_2'
}

\sEntry{Binary-operator-reduce}{A binary operator applied to two values is
replaced by the result of the operation.
Integers are arbitrary-precision; if one of the operands is a float, the
other operand is converted to a float, and if one of them is complex, the
other is converted to complex (see the numbers section of this specification).
Operands of other types, for example a boolean operand
of an arithmetic operator, cause a \texttt{TypeError}, and division by zero
causes a \texttt{ZeroDivisionError}; the reduction is then stuck.}
\gentzen{
  v_1\ \sVar{op}\ v_2 \mbox{ evaluates to } v
}{
  v_1\  \sVar{binary-operator} \ v_2
  \ \Rightarrow \ 
  v
}

\sEntry{And-shortcut}{An expression with binary operator
\sCode{and} whose left sub-expression is \sCode{False} is reduced to \sCode{False},
without evaluating the right sub-expression.}
\gentzen{
}{
  \sCode{False and }e \ \Rightarrow \ \sCode{False}
}

\sEntry{And-continue}{An expression with binary operator
\sCode{and} whose left sub-expression is \sCode{True} is reduced to
the right sub-expression.}
\gentzen{
}{
  \sCode{True and }e \ \Rightarrow \ e
}

\sEntry{Or-shortcut}{An expression with binary operator
\sCode{or} whose left sub-expression is \sCode{True} is reduced to
\sCode{True}, without evaluating the right sub-expression.}
\gentzen{
}{
  \sCode{True or }e \ \Rightarrow \ \sCode{True}
}

\sEntry{Or-continue}{An expression with binary operator
\sCode{or} whose left sub-expression is \sCode{False} is reduced to
the right sub-expression.}
\gentzen{
}{
  \sCode{False or }e \ \Rightarrow \ e
}

The left operand of \sCode{and} and \sCode{or} is evaluated first,
as in the rules for the other binary operators (\textbf{Left-binary-reduce}).
A left operand that is a value other than \sCode{True} or \sCode{False} causes
a \texttt{TypeError}.

\subsection*{Expressions: Unary operators}

\sEntry{Unary-reduce}{An expression with a unary operator
can be reduced if its operand can be reduced.}
\gentzen{
  e \ \Rightarrow \ e'
}{
  \sVar{unary-operator}\ e
  \ \Rightarrow \ 
  \sVar{unary-operator}\ e'
}

\sEntry{Unary-operator-reduce}{A unary operator (\sCode{-}, \sCode{+} or
\sCode{not}) applied to a value is replaced by the result of the operation.
\sCode{not} requires a boolean operand.}
\gentzen{
  \sVar{op}\ v \mbox{ evaluates to } v'
}{
  \sVar{unary-operator}\ v
  \ \Rightarrow \ 
  v'
}

\subsection*{Expressions: Conditional expressions}

\sEntry{Conditional-expression-predicate}{A conditional expression is
reducible if its predicate is reducible.}
\gentzen{
  e_1\ \Rightarrow\ e_1'
}{
  e_2\ \sCode{ if }\ e_1\ \sCode{ else }\ e_3
  \ \Rightarrow \ 
  e_2\ \sCode{ if }\ e_1'\ \sCode{ else }\ e_3
}

\sEntry{Conditional-expression-consequent}{A conditional expression whose
   predicate is \sCode{True} reduces to its consequent.}
\gentzen{
}{
  e_2\ \sCode{ if True else }\ e_3 \ \Rightarrow \ e_2
}

\sEntry{Conditional-expression-alternative}{A conditional expression whose
   predicate is \sCode{False} reduces to its alternative.}
\gentzen{
}{
  e_2\ \sCode{ if False else }\ e_3 \ \Rightarrow \ e_3
}

\subsection*{Expressions: Function application}

\sEntry{Application-functor-reduce}{A function application
can be reduced if its functor expression can be reduced.}
\gentzen{
  e \ \Rightarrow \ e'
}{
  e\  \sCode{(}\ \sVar{expressions} \ \sCode{)}
  \ \Rightarrow \ 
  e'\  \sCode{(}\ \sVar{expressions} \ \sCode{)}
}

\sEntry{Application-argument-reduce}{A function application
can be reduced if one of its argument expressions can be reduced and all
preceding arguments are values.}
\gentzen{
  e \ \Rightarrow \ e'
}{
  v\  \sCode{(}\ v_1 \ldots v_i \sCode{, } e \sCode{, }\ldots\ \sCode{)}
  \ \Rightarrow \ 
  v\  \sCode{(}\ v_1 \ldots v_i \sCode{, } e' \sCode{, }\ldots\ \sCode{)}
}

\sEntry{Function-definition-application-reduce}{The application of a function
definition, to as many arguments as it has parameters, all of which are values,
is replaced by its body in which the parameters are replaced by the arguments,
and the name of the function is replaced by the function itself (so that
recursive calls can be applied). If the body is a suite that does not start with
a \sCode{return}, the application is replaced by the block expression of
the suite.}
\gentzen{
  f = \sCode{def }n\sCode{(}x_1\sCode{, }\ldots\sCode{, }x_k\sCode{):}\ b
}{
  f\ \sCode{(}\ v_1 \sCode{, }\ldots\sCode{, }v_k\ \sCode{)}
  \ \Rightarrow \ 
  b [x_1 := v_1]\ldots[x_k := v_k]
  [n := f]
}

\sEntry{Lambda-application-reduce}{The application of a lambda
expression, to as many arguments as it has parameters, all of which are values,
is replaced by its body in which the parameters are replaced by the arguments.
A lambda expression that was assigned to a name $n$ (\sCode{n = lambda ...})
is applied in the same way as a function definition, including the
substitution $[n := f]$.}
\gentzen{
  f = \sCode{lambda }x_1\sCode{, }\ldots\sCode{, }x_k\sCode{: }e
}{
  f\ \sCode{(}\ v_1 \sCode{, }\ldots\sCode{, }v_k\ \sCode{)}
  \ \Rightarrow \ 
  e [x_1 := v_1]\ldots[x_k := v_k]
}

If the number of arguments is different from the number of parameters,
the reduction is stuck with a \texttt{TypeError}. The same holds for the
application of a value that is not a function.

\sEntry{Builtin-application-reduce}{The application of a built-in function
(for example \sCode{abs}, \sCode{math\_sqrt}, \sCode{pair}, \sCode{head},
\sCode{tail} and the other functions of the linked list library), to arguments
that are all values, is replaced by the result of the function in one step.
The functions \sCode{print} and \sCode{print\_llist} additionally write
to the output, which the stepper shows from the step after this contraction
onwards.}
\gentzen{
  f \mbox{ is a built-in function}\quad f(v_1,\ldots,v_k) = v
}{
  f\ \sCode{(}\ v_1 \sCode{, }\ldots\sCode{, }v_k\ \sCode{)}
  \ \Rightarrow \ 
  v
}

\pagebreak

\section*{Substitution}

Substitution $p[x := v]$ replaces the free occurrences of $x$ in $p$ by
$v$. It is defined by the following cases; for the expressions that are not
listed (operators, applications, pairs, conditional expressions), the
substitution is applied to every sub-expression.

\sEntry{Identifier}{An identifier with the same name as $x$ is substituted
with $v$; other identifiers are left alone.}
\gentzen{}{
  x[x := v] \ = \ v \qquad\qquad
  \sVar{name}[x := v] \ = \ \sVar{name}\quad (\sVar{name} \neq x)
}

\sEntry{Compound expressions}{All occurrences of $x$ in the
sub-expressions are substituted with $v$.}
\gentzen{}{
  (e_1 \ \sVar{binary-operator} \ e_2)[x := v] \ = \ 
  e_1[x := v] \ \sVar{binary-operator} \ e_2[x := v]
}

\sEntry{Lambda expression}{If $x$ is one of the parameters, the
lambda expression is left alone. Otherwise the body is substituted. If a
parameter of the lambda expression occurs free in $v$, the parameter is
renamed first, to avoid that the parameter captures that occurrence.}
\gentzen{
  x \notin \{x_1, \ldots, x_k\}\quad
  \{x_1, \ldots, x_k\} \cap \mathit{free}(v) = \emptyset
}{
  (\sCode{lambda }x_1\sCode{, }\ldots\sCode{, }x_k\sCode{: }e)[x := v]
  \ = \ 
  \sCode{lambda }x_1\sCode{, }\ldots\sCode{, }x_k\sCode{: }e[x := v]
}

\sEntry{Function definition}{If $x$ is the name of the function, or one of the
parameters, the definition is left alone. Otherwise, the body is substituted,
after the renaming of parameters and local names that occur free in $v$, as for
lambda expressions. If $x$ is assigned to in the body of the definition, $x$ is
local to the function (as in Python): $v$ is not substituted into the
body, and reading $x$ before the assignment is an
\texttt{UnboundLocalError}.}
\gentzen{
  x \notin \{n, x_1, \ldots, x_k\}\quad x \mbox{ is not assigned in } b
}{
  (\sCode{def }n\sCode{(}x_1\sCode{, }\ldots\sCode{, }x_k\sCode{):}\ b)[x := v]
  \ = \ 
  \sCode{def }n\sCode{(}x_1\sCode{, }\ldots\sCode{, }x_k\sCode{):}\ b[x := v]
}

\sEntry{Statement sequences}{In a sequence of statements, the statements are
substituted from left to right. After a statement that declares $x$ (an
assignment to $x$, or a definition of $x$), the remaining statements are left
alone, because they refer to that declaration.}
\gentzen{
  \mbox{the first statement of } s \mbox{ declares } x
}{
  (\sVar{statement}_1\ \sVar{statement}_2 \ldots)[x := v]
  \ = \ 
  \sVar{statement}_1[x := v]\ \sVar{statement}_2\ldots
}

\sEntry{Assignment}{In an assignment, only the right-hand side is
substituted.}
\gentzen{}{
  (\sVar{name}\sCode{ = }e)[x := v] \ = \ \sVar{name}\sCode{ = }e[x := v]
}

\section*{Free names}

The set $\mathit{free}(p)$ of free names of a program, statement or expression
$p$ is defined as usual: an identifier is free; a lambda expression binds
its parameters; a function definition binds its name, its parameters and the
names that are assigned to in its body; and, in a sequence of statements, an
assignment or function definition binds its name in the statements that
follow it. Names of built-in functions are not free: they are never
substituted.

\section*{Steps and their limit}

The stepper presents an evaluation as a sequence of \emph{steps}, numbered
from 0 on the step slider:
\begin{itemize}
\item Step 0 shows the program before evaluation (``Start of evaluation'').
\item Each contraction contributes two steps: the program with the redex
  highlighted (the explanation says what is about to happen, for example
  ``Evaluating binary expression''), followed by the program after the
  contraction with the result highlighted (``Evaluated binary expression
  \ldots'').
  The program after a contraction is identical to the program before the next
  contraction; only the highlighting differs. If a contraction removes a
  statement without leaving a result in its place (for example
  \textbf{Eliminate-assignment}), nothing is highlighted after it.
\item The last step is ``Evaluation complete'', ``Evaluation stuck'' (in
  which case the previous step explains the error), or ``Maximum number
  of steps exceeded''.
\end{itemize}

The number of the last step does not exceed the \emph{step limit} that the
user can configure (default: 1000). A limit of $i$ allows
$\max(0, \lfloor (i - 1) / 2 \rfloor)$ contractions; the last step is
numbered $2\cdot\lfloor (i - 1) / 2 \rfloor + 1$ if the evaluation
completes within the limit. The same limit applies to the environment stepper
and to the CSE machine, whose step slider also counts from 0.
